\section*{Bab \ref{ch:b3:quantum-statistics} --- Statistika Kuantum}

\begin{solution}{exo:b3:quantum-statistics:1}
(a) $\qty{7.1}{eV}$, $\qty{8.2e4}{K}$. (b) $\qty{3.2}{eV}$,
$\qty{3.7e4}{K}$. (c) Dengan $m = 3u$: $E_{\text{F}} \approx
\qty{4e-4}{eV}$, $T_{\text{F}} \approx \qty{5}{K}$ (yakni nilai gas
idealnya; interaksinya menurunkannya). (d) Elektron logam: merosot pada
tiap suhu di Bumi; $^3$He: hanya di bawah beberapa kelvin;
$T = \qty{300}{K}$ tak pernah menyentuh skala Fermi sebuah logam.
\end{solution}

\begin{solution}{exo:b3:quantum-statistics:2}
(a) $\int_0^{E_{\text{F}}}\epsilon\,g\,\dd\epsilon/\int g\,\dd
\epsilon = \tfrac35 E_{\text{F}}$ dari $g \propto \sqrt\epsilon$.
(b) $v_{\text{F}} = \sqrt{2E_{\text{F}}/m_{\text{e}}} =
\qty{1.6e6}{m/s}$; akar rata-rata kuadratnya $\sqrt{3/5}\,v_{\text{F}} \approx
\qty{1.2e6}{m/s}$. (c) $\tfrac32 k_{\text{B}}T = \tfrac35
E_{\text{F}}$: $T \approx \qty{3.3e4}{K}$. (d) Tak ada tempat untuk
pergi: tiap keadaan yang lebih rendah sudah terisi, sehingga tak ada foton
yang dapat dipancarkan dan tak ada tumbukan yang dapat memperlambatnya
--- laju tanpa kalor, berkat pengecualian.
\end{solution}

\begin{solution}{exo:b3:quantum-statistics:3}
(a) $P = 0.4 \times \num{8.5e28} \times \qty{7.1}{eV} =
\qty{3.9e10}{Pa} \approx \num{4e5}$ atmosfer. (b) Tarikan Coulomb
kisi ionnya: logamnya adalah gencatan senjata antara
keduanya (\cref{exo:b3:quantum-statistics:11}). (c) Pemampatan menaikkan
$E_{\text{F}}$ secara curam ($n^{2/3}$): logam melawan. (d) $K \sim
\tfrac53 P \approx \qty{65}{GPa}$ berbanding nilai terukurnya
\qty{140}{GPa}: separuh kekakuan tembaga adalah pengecualian murni.
\end{solution}

\begin{solution}{exo:b3:quantum-statistics:4}
(a) $T_{\text{c}} \approx \qty{0.4}{\micro K}$ --- kondensat JILA
muncul pada \qty{170}{nK}: fisika yang sama, dengan kerapatan puncak
yang sedikit lebih rendah. (b) $T_{\text{c}} \propto 1/m$: $\times87 \approx
\qty{35}{\micro K}$. (c) Pada kerapatan tinggi dan dingin sedemikian,
tiap gas kecuali hidrogen yang terkutub spinnya \emph{memadat};
tumbukan tiga benda membuat molekul. Kiatnya adalah gas encer yang
lengai dan cukup dingin sehingga $\lambda_T$ mengimbangi keencerannya. (d)
$\lambda_{T_{\text{c}}} = h/\sqrt{2\pi mk_{\text{B}}T_{\text{c}}}
\approx \qty{0.3}{\micro m}$: $n\lambda^3 = \num{e20} \times
\num{2.7e-20} \approx 2.7$, sebagaimana dituntut aturannya.
\end{solution}

\begin{solution}{exo:b3:quantum-statistics:5}
(a) $\gamma T = \beta T^3$ pada $T^2 = \gamma/\beta$: untuk tembaga
$T \approx \qty{3}{K}$. (b) $C_{\text{el}}/C \approx
(\pi^2/2)(300/\num{8.2e4})/3 \approx 0.6\%$. (c) Pada \qty{0.1}{K}
suku kisinya sudah runtuh sebagai $T^3$: elektronnya mengangkut
$\sim10^3$ kali lebih banyak --- kalorimetri menjadi spektroskopi
elektron. (d) $C/T = \gamma + \beta T^2$ adalah garis lurus dalam
$T^2$: titik potong $\gamma$ mengukur \omterm{prop:b3:schrodinger-three-dimensions:counting}{kerapatan keadaan} pada
aras Ferminya, kemiringan $\beta$ mengukur suhu Debyenya --- dua bilangan
tiap logamnya dari satu plot.
\end{solution}

\begin{solution}{exo:b3:quantum-statistics:6}
(a) $N_{\text{e}} = M/2m_{\text{n}} = \num{6e56}$; $n =
\num{5.5e35}\,\unit{m^{-3}}$; $E_{\text{F}} \approx
\qty{0.24}{MeV}$. (b) Separuh \omterm{def:b3:relativistic-dynamics:momentum-energy}{energi diam} elektronnya (yakni tepi
rezim relativistiknya), namun tiga ratus kali $k_{\text{B}}T$
(\qty{0.9}{keV}): membara putih dan dingin secara kuantum sekaligus. (c)
$P_{\text{deg}} \approx \qty{9e21}{Pa}$ berbanding $GM^2/R^4 \sim
\num{e22}\,\unit{Pa}$: seorde --- bintangnya persis kesetimbangan
ini. (d) Massa yang lebih besar: sisi gravitasinya tumbuh sebagai $M^2$, dan
bintangnya harus menyusut ($R \propto M^{-1/3}$), lalu menyeret $E_{\text{F}}$
menuju pelunakan relativistiknya.
\end{solution}

\begin{solution}{exo:b3:quantum-statistics:7}
(a) $\rho \approx \qty{5e17}{kg/m^3}$: kira-kira dua kali kerapatan
inti --- yakni sebuah inti seukuran kota. (b) $n_{\text{n}}
\approx \qty{3e44}{m^{-3}}$: $E_{\text{F}} \approx \qty{90}{MeV}$,
yakni sepersepuluh $m_{\text{n}}c^2$: kenisbian sedang mengetuk. (c) \omterm{thm:b3:quantum-statistics:fermi}{Energi
Fermi} elektronnya menaik melewati ongkos \qty{0.78}{MeV} bagi
$\text{e} + \text{p} \to \text{n} + \nu$: penangkapannya menjadi
menguntungkan, lalu bintangnya menukarkan elektron dan protonnya dengan
neutron. (d) $\qty{5}{cm^3} \times \qty{5e17}{kg/m^3} =
\qty{2.5e12}{kg}$: satu sendok teh yang lebih berat daripada pegunungan.
\end{solution}

\begin{solution}{exo:b3:quantum-statistics:8}
(a) Mendinginkan sebuah rongga sekadar \emph{menyingkirkan} fotonnya
(dindingnya menyerapnya): tanpa cacah yang harus kekal, $\mu$ tetap
terpancang pada nol dan tak pernah ada kelebihan yang menumpuk ke dalam
ragam dasarnya --- gasnya meredup dan bukan memampat. (b) Cacah atomnya
kekal pada tiap skala waktu laboratorium. (c) Zat warnanya berulang kali
menyerap lalu memancarkan kembali fotonnya \emph{tanpa kehilangan
fotonnya} pada skala waktu penermalannya: yakni cacah foton yang praktis
kekal (dan massa efektif dari rongganya) --- lalu pemampatannya pun
muncul.
(d) Fonon: tercipta dan termusnah dengan bebas, $\mu = 0$: tanpa
pemampatan --- ``gas fonon'' sebuah kristal sekadar membeku.
\end{solution}

\begin{solution}{exo:b3:quantum-statistics:9}
(a) $T_{\text{c}}^{\text{ideal}} \approx \qty{3.1}{K}$ berbanding
nilai terukurnya $T_\lambda = \qty{2.17}{K}$. (b) Helium cair itu rapat
dan berinteraksi kuat: ``ideal'' membuang interaksi yang
menggeser (dan membentuk ulang) peralihannya. (c) Atom $^3$He bersifat fermion:
tanpa pemampatan tanpa pemasangan; interaksi yang menarik mengikat
pasangan bergaya Cooper hanya di bawah \qty{2.6}{mK}, lalu cairan
berpasangan itu mengalir adiluncur. (d) Sembilan fermion: $^6$Li adalah fermion (sedangkan
$^7$Li adalah boson) --- pasangan isotop litium ini memungkinkan satu
laboratorium mempelajari laut Fermi dan sebuah kondensat dalam satu oven.
\end{solution}

\begin{solution}{exo:b3:quantum-statistics:10}
(a) Gantikanlah $g \propto \sqrt\epsilon$ lalu keluarkanlah $\beta$ lewat penskalaan.
(b) Nilai integralnya memberikan $N_{\text{exc}} =
2.612\,V/\lambda_T^3$; menetapkan $N_{\text{exc}} = N$ mendefinisikan
$T_{\text{c}}$. (c) Di bawah $T_{\text{c}}$, $N_{\text{exc}}(T) =
N(T/T_{\text{c}})^{3/2}$; sisanya adalah $N_0$. (d) Dalam dua
dimensi $g(\epsilon)$ bersifat tetap dan $\int\dd\epsilon/(\eu^{
\beta\epsilon} - 1)$ menyebar di dasarnya: keadaan tereksitasinya
selalu dapat menyerap semuanya --- tanpa populasi keadaan dasar yang
makroskopis pada $T > 0$ mana pun.
\end{solution}

\begin{solution}{exo:b3:quantum-statistics:11}
(a) $\dd\epsilon/\dd n = \tfrac23 An^{-1/3} - \tfrac13 Bn^{-2/3} =
0$ pada $n^{1/3} = B/2A$: yakni minimum yang sejati. (b) Dengan $A \sim
\hbar^2/m_{\text{e}}$ dan $B \sim e^2/4\pi\varepsilon_0$, jarak
kesetimbangannya $n^{-1/3} \sim \hbar^2 4\pi\varepsilon_0/
m_{\text{e}}e^2 = a_0$: logam bersifat rapat karena \omterm{thm:b3:hydrogen-atom:levels}{jejari Bohr}
mengatakannya demikian. (c) $\epsilon_{\min} \sim -B^2/A \sim -E_{\text{I}}$
--- yakni kelekatan berorde elektronvolt, sebagaimana terukur. (d) \omterm{thm:b3:quantum-statistics:fermi}{Tekanan
kemerosotan} adalah separuh penolak dalam gencatan senjata itu: tanpanya
suku Coulombnya akan meremukkan kisinya menjadi satu titik --- limbak
materi adalah milik Pauli.
\end{solution}

\begin{solution}{exo:b3:quantum-statistics:12}
(a) Elektron aktifnya: $g(E_{\text{F}})k_{\text{B}}T$; energi
masing-masing: $\sim k_{\text{B}}T$; lalu turunkanlah $E_{\text{tambahan}} \sim
g k_{\text{B}}^2T^2$. (b) $g(E_{\text{F}}) = 3N/2E_{\text{F}}$
memberikan $C \sim 3Nk_{\text{B}}T/T_{\text{F}}$ --- nilai eksaknya
$\pi^2/2$ menggantikan 3. (c) Spin yang dapat berbalik $\propto T$; tiap
spinnya menyumbang penyearahan $\propto \mu^2B/k_{\text{B}}T$: $T$-nya
saling meniadakan --- kerentanan Pauli bersifat datar sedangkan kerentanan
Curie (\cref{exo:b3:canonical-ensemble:11}) turun sebagai $1/T$: yakni
tanda pengenal kemerosotan. (d) Fermion yang merosot hanya menanggapi di
permukaannya: kalikanlah jawaban klasik mana pun dengan $\sim
T/T_{\text{F}}$, atau batasilah ia pada kulit Ferminya.
\end{solution}

\begin{solution}{pb:b3:quantum-statistics:1}
\textbf{1.} Energi yang terlepas ketika bintangnya dirakit dari materi
yang tersebar --- bernilai negatif, dan makin dalam untuk yang makin ringkas.
\textbf{2.} $E_{\text{F}} \sim (\hbar^2/m_{\text{e}})
(N_{\text{e}}/R^3)^{2/3}$.
\textbf{3.} $E_{\text{k}} \sim N_{\text{e}}E_{\text{F}} \sim
\hbar^2N_{\text{e}}^{5/3}/m_{\text{e}}R^2$.
\textbf{4.} Meminimumkan $A/R^2 - B/R$ memberikan $R = 2A/B$: $R \sim
\hbar^2N_{\text{e}}^{-1/3}/(Gm_{\text{e}}m_{\text{n}}^2\mu_e^2)$.
\textbf{5.} $N_{\text{e}} \propto M$: $R \propto M^{-1/3}$ ---
menambahkan massa justru \emph{menyusutkan} bintangnya, sehingga
kerapatan dan $E_{\text{F}}$
tumbuh tanpa batas: penopang taknisbiannya hidup dari utang.
\textbf{6.} $N_{\text{e}} = \num{6e56}$: $R \sim \qty{2000}{km}$
--- orde yang tepat di samping \qty{5800}{km} Sirius B (penskalaan kami
membuang faktor berorde beberapa satuan).
\textbf{7.} $E_{\text{F}} \sim m_{\text{e}}c^2$ pada $n \sim
(m_{\text{e}}c/\hbar)^3 \approx \qty{2e37}{m^{-3}}$, yakni $\rho
\sim \num{e10}$--$\num{e11}\,\unit{kg/m^3}$: yakni rezim katai
terberat.
\textbf{8.} Tiap elektronnya: $\epsilon \sim \hbar c\,n^{1/3}$;
seluruhnya: $\hbar cN_{\text{e}}^{4/3}/R$.
\textbf{9.} Dengan kedua sukunya $\propto 1/R$, jejarinya lenyap:
kemantapannya ditentukan oleh apakah koefisien kinetiknya melampaui
koefisien gravitasinya --- yakni kriteria massa yang murni.
\textbf{10.} Samakanlah: $N_{\text{e,c}} \sim (\hbar c/
Gm_{\text{n}}^2\mu_e^2)^{3/2}$, dan $M_{\text{Ch}} =
\mu_em_{\text{n}}N_{\text{e,c}}$: diperoleh rumus yang dinyatakan itu.
\textbf{11.} $\hbar c/Gm_{\text{n}}^2 = \num{1.7e38}$;
$M_{\text{Ch}} \sim (\num{1.7e38})^{3/2}m_{\text{n}}/4 \approx
\qty{9e29}{kg} \approx 0.5\,M_\odot$ --- perhitungan eksaknya
menajamkan orde besarnya menjadi $1.4\,M_\odot$.
\textbf{12.} $\hbar$ (tekanannya bersifat kuantum), $c$ (yakni
pelunakan relativistiknya), $G$ (yakni lawannya): sebuah bilangan
seukuran bintang yang dibangun dari tiga tetapan mikroskopis --- yakni
mekanika kuantum yang membuat undang-undang bagi benda berisi $10^{57}$ partikel.
\textbf{13.} Kemerosotan neutron, pada jejari yang lebih kecil sebesar
$\sim m_{\text{e}}/m_{\text{n}}$ (dikalikan pembukuan $\mu_e$-nya):
yakni kilometer --- inilah bintang neutron.
\textbf{14.} Ketiga tetapan yang sama dengan $m_{\text{n}}$ di
mana-mana memberikan skala massa yang sama: bintang neutron berhenti di
dekat $2\,M_\odot$; di baliknya, tak ada hukum tekanan yang menang ---
yakni lubang hitam.
\textbf{15.} $GM^2/R \approx \qty{5e46}{J}$: beberapa ratus kali
seluruh keluaran Matahari selama sepuluh miliar tahun, yang terlepas
dalam hitungan detik.
\textbf{16.} Supernova 1987A: dua lusin neutrino, yang tiba beberapa
jam sebelum cahayanya memijar, membawa persis energi yang diramalkan
--- yakni keruntuhan yang terdengar langsung.
\textbf{17.} Keruntuhan teras: teras besi sebuah bintang masif meledak
ke dalam melewati tiap lantai tekanannya; selubungnya tertiup pergi.
Jenis Ia: sebuah \omterm{ex:b3:quantum-statistics:dwarf}{katai putih} yang disuapi sampai ambangnya menyalakan
karbonnya dalam kilatan termonuklir yang melepas ikatan seluruh bintangnya.
\textbf{18.} Massa pemicu yang sama, bahan bakar yang sama, pelepasan
energi yang sama: kurva cahaya yang nyaris seragam --- yakni bom yang
terkalibrasi.
\textbf{19.} Dengan mengetahui kecerlangan mutlaknya, kecerlangan
teramatinya memberikan jaraknya lewat kuadrat terbalik; tanpa sifat
``baku'' tangganya runtuh.
\textbf{20.} Terlalu redup berarti terlalu jauh: pemuaiannya ternyata
sedang memecut --- yakni energi gelap, yang buktinya bersandar pada
keseragaman lilinnya (karena itulah fisikanya diurus dengan hati-hati).
\textbf{21.} Asas pengecualian $\to$ \omterm{thm:b3:quantum-statistics:fermi}{tekanan kemerosotan} $\to$ sebuah
massa batas yang semesta $\to$ ledakan yang dibakukan $\to$ pemecutan
alam semesta yang terukur.
\textbf{22.} Ikutilah matematikanya ke mana pun ia menuju: perilaku
``mustahil'' bintang di balik batasnya ternyata adalah bintang neutron
dan lubang hitam --- keduanya kini terkatalog ribuan.
\textbf{23.} Pada 1915 tak ada fisika yang dikenal mengizinkan materi
satu ton tiap sentimeter kubik; pada 1926 statistika Fermi--Dirac
menjadikannya keadaan baku tiap gas yang dingin dan rapat: pengamatannya
menunggu statistikanya, bukan sebaliknya.
\textbf{24.} Bahwa katai berada secara mantap sepanjang suatu rentang
massa (yakni cabang $R(M)$-nya), dan bahwa cabang itu \emph{berakhir}
pada $1.4\,M_\odot$: keduanya tergambar dalam histogram itu.
\textbf{25.} $(\hbar c/Gm_{\text{n}}^2)^{3/2}m_{\text{n}} \approx
1.4\,M_\odot$: tiga tetapan menetapkan \omterm{ex:b3:quantum-statistics:dwarf}{katai putih} terberat ---
sehingga menetapkan pula massa pemicu lilin baku yang menimbang
nasib alam semesta.
\end{solution}
